\section{Gate instruction grammar and arity}

\subsection{General gate shape}

\begin{lstlisting}
GATE -I input_or_controls -O target
\end{lstlisting}
The uppercase word selects a gate. \code{-I} lists the incoming wire for a
one-qubit operation or the controls for a controlled operation. \code{-O}
names the wire that is mutated. Except for SWAP, gate instructions do not
create a separate output wire automatically.

\begin{lstlisting}
H -I Q -O Q
CNOT -I Control -O Target
\end{lstlisting}
In the first row, Q is listed twice because the state flows through H on one
wire. In the second, Control is preserved and Target is conditionally changed.

\subsection{Why arity matters}

\begin{center}
\begin{tabular}{lll}
\toprule
Gate family & Required \code{-I} count & Required \code{-O} count\\
\midrule
X, Y, Z, H, S, T, PHASE, NOT & 1 & 1\\
CNOT, CZ, CY & 1 & 1\\
CCNOT, AND, NAND, OR, XOR & 2 & 1\\
SWAP & 2 & 2\\
MEASURE & 0 or 1 & 0\\
\bottomrule
\end{tabular}
\end{center}

All counts are exact except \code{MEASURE}. For ordinary single-qubit gates,
name the same wire in \code{-I} and \code{-O}. The output is the mutated target;
for controlled gates, \code{-I} contains controls.

\subsection{Validation and common mistakes}

Arity is part of the language contract. A CNOT with two inputs is not silently
converted into CCNOT, and a CCNOT with one input is not converted into CNOT.
The compiler reports whether an instruction has too few or too many inputs or
outputs.

\begin{lstlisting}
# Wrong: H needs one output.
H -I Q

# Wrong: CNOT needs exactly one control.
CNOT -I A B -O Target

# Right:
CCNOT -I A B -O Target
\end{lstlisting}

The target is removed from the compiler's control list if the same resolved
wire appears on both sides. This is necessary for the normal one-qubit form.
For controlled gates, do not also list the target as an input; doing so either
violates arity or obscures the intended control relationship.

